\section{人类的位置：能效锚点与异质性阻尼}

智能的模态分化把任务拆成探索、控制与验证，也使人的位置需要按任务重新核算。这里考察三个方面：完成同一任务的系统能耗，多个回路共同作用时的稳定性，以及结果进入行动的准入条件。前两项需要工程比较，第三项还涉及风险分配和社会授权。人的能力、机器的能力与制度安排共同决定具体分工。

\subsection{湿件的模拟计算优势：能量账本上的位置}

把人放进模态分化里，先要换一个比较口径。按能力条目逐个打分的方式每一轮都会变，因为机器的成本曲线还在下降；相对稳定的口径是能量账本：完成同一件工作，两边各要支付多少功率、多少散热、多少自重搬运。这个口径不新，但在讨论智能时经常被跳过，因为智能看起来像是纯信息活动，纯信息活动不消耗功率。第2章已经把这个印象处理掉了，兰道尔原理给出的下界虽然小到可以忽略，信息处理仍然需要一个物理基座，而这个基座要供电、要冷却、要维护、要更新。\cite{landauer1961} 一旦把基座算进来，能量账本上的对比就不再是“人脑对芯片”，而是“整套生物系统对整套机器系统”。

能耗比较必须采用相同的系统边界。人脑持续功率约在20瓦量级，人体整体代谢则随活动强度变化。\cite{smil2015power} 机器需要计入感知、计算、执行与辅助设施，人也需要计入维持身体和完成任务的代谢支出。把两侧都按互不重叠的功能项目核算，可以写成

\begin{equation}
E_s=E_{s,\text{感知与计算}}+E_{s,\text{执行}}+E_{s,\text{辅助}},\qquad s\in\{\text{人},\text{机器}\}
\end{equation}

各项单位均为焦耳，核算相同任务、质量标准和完成时间。执行项包括对外做功与移动自身，辅助项包括维持运行所需的其他支出，各项不得重复。人体各功能共享组织与代谢通路，难以逐项测量时，可以使用相同边界内的总能耗；机器也应采用相应总量。比较结果由测量给出，分项多少本身不构成能效优势。

身体结构会改变完成任务所需的控制。抓取时，肌腱、皮肤和关节的顺应性可以吸收一部分接触误差，神经控制与被动力学共同完成动作。机器也可以利用柔性结构、顺应控制和形态计算，将部分修正交给身体。两类系统的成本取决于具体实现：哪些偏差由结构吸收，哪些需要传感、计算和主动执行，哪些会转成维护与安全冗余。

配送、护理、维修、装卸和巡检包含大量非标准现场。门牌、楼梯、障碍、天气和他人行为都在变化，任务还可能需要沟通与临机处置。机器进入这些环境，要计算感知冗余、远程接管与现场支持；人的工作同样包含培训、疲劳、受伤和组织支持。\cite{moravec1988,pfeiferbongard2006} 能耗只是比较的一项，可靠性、服务质量和维护成本也会改变可行的分工。

标准货架、固定航线和受控产线，能够减少感知与动作中的不确定性，使自动化更容易形成稳定收益。开放环境则需要更多在线修正。人的身体已经积累了演化与学习形成的适应能力，机器的结构与训练也在改变。哪一侧占优，需要在相同任务和质量要求下比较，不能从环境名称直接读出。

\subsection{系统的异质性阻尼}

第二个问题是共同响应。多个参与者若在同一信号下同时行动，局部修正可能叠成一次集体冲击。不同的响应节奏可以分散这种冲击，也可能增加协调困难。判断异质性是否有利，需要考察信号、连接结构与反馈方向。

先看一个简化情形：若若干线性通道接受同一频率的输入，其输出可以用复数振幅相加。设第$i$个通道的幅值增益为$g_i$、相位为$\varphi_i$，合成幅值为

\begin{equation}
G_{\text{总}}=\Bigl|\sum_{i=1}^{N} g_i\,e^{\mathrm{i}\varphi_i}\Bigr|
\end{equation}

$g_i$与$G_{\text{总}}$无量纲，$\varphi_i$以弧度计。同相时，合成幅值达到各项幅值之和；相位分散时，部分响应可能抵消。这个关系描述特定频率上的输出叠加。闭环是否稳定，还取决于完整的环路传递函数与反馈连接，合成幅值越过1本身不能判定自激振荡。

退避、随机抖动、分批执行和独立复核，可以减少同一信号引发的同步行动。人的参与有时带来不同信息与判断路径，疲劳和延迟也可能增加失误。需要保留的是可以被验证的分散与纠错功能，具体由人工、软件或制度提供，要按结果比较。

异质性与协调需要共同设计。执行节奏过于一致可能放大冲击，信息脱节则可能使行动相互抵消。自动化中的人工监督也需要可获得的信息、足够的反应时间与持续训练。\cite{bainbridge1983ironies,parasuraman1997humans} 人是否处在有用的反馈位置，由这些条件决定。

\subsection{意义作为停机条件}

第三个问题是何时结束搜索、允许结果进入行动。算法可以预设时间预算、误差阈值或收敛条件，自动决定停止；社会还需要确定怎样的结果足以被依赖。搜索的终止与行动的准入在这里分开。

优化目标可以有界，也可以无界；多目标问题可能形成一组互不支配的候选，仍可通过明确规则停止计算。部署需要的是一组经过说明的容许条件。设$\mathcal A$为满足结果质量与风险要求的状态集合，搜索首次进入该集合的时刻记为

\begin{equation}
t^{\ast}=\inf\bigl\{t\ge0:\ \mathbf{x}_t\in\mathcal{A}\bigr\},\qquad \inf\varnothing=+\infty
\end{equation}

$\mathbf x_t$为搜索状态，$t^*$以时间计。搜索始终未进入$\mathcal A$时，首次到达时刻为无穷，实际程序还可以按预算终止并报告失败。$\mathcal A$既可由规则自动检查，也可包含需要进一步试验或授权的要求。这个表达式记录准入条件与到达过程，条件的内容须另行确定。

人的判断主要进入条件的制定、修订与争议处理。哪些错误可接受，哪些风险需要继续验证，谁可以要求暂停，涉及承受者、部署者和公共制度。模型可以参与拟定和检验这些条件，采用什么标准则需要相应授权。指标进入奖惩以后还会改变行为，因此准入标准也需要申诉、复核与修订机制。\cite{goodhart1975problems,campbell1979assessing}

授权与专业能力须同时考察。承担风险的一方应当能够了解标准、表达异议并获得救济；负责审核的一方则需要证据和能力。人工判断可能失误，自动检查也可能遗漏，制度需要让错误可发现、决定可追溯、责任可落实。第6章将在价值排序和公共裁决中继续处理这组关系。

这三个方面分别给出能耗、稳定性和准入的考察方法。能耗优势要由同任务比较确定，稳定作用要由反馈结构检验，决定资格要由权利与授权安排确立。人的位置随技术和制度共同变化。以这些条件评价具体分工，比逐项宣布某种能力属于谁，更接近现实系统的运行方式。

分工回答的是“谁来承担”，人类的位置回答的是“人站在哪一项上”，两者都没有回答“在什么地方承担得起”。这三种模态最终都要落到具体的物理系统上，而现实系统并不按智能的分工方式划分自己：一条道路、一张电网、一套时刻表，同时承担着资源账本、时间刚性和结构耦合，它们不会因为智能分工合理而自动变轻，也不会因为某一环节被局部优化而整体降低黏滞。把一个环节做得更聪明很容易，让一整套位移系统变得更轻很难，而这类系统恰好最擅长把单点优化的收益自己吞掉。
